% !TEX root = ../Main.tex
\chapter{力的合成（特殊三角形）}

两个等大的力 $F_1=F_2=F$ 合成时，遇到几个特殊夹角会得到很整齐的合力。下面把三种特殊夹角（$60^\circ$、$90^\circ$、$120^\circ$）的受力图画出来，合力用平行四边形的对角线表示，请按平行四边形法则求出各自的合力大小。

\section{三种特殊夹角}

\begin{center}
\begin{tikzpicture}[>=Stealth, scale=1]
    % 60°
    \begin{scope}
        \draw[->, thick] (0,0) -- (60:1.8) node[above] {\small $F$};
        \draw[->, thick] (0,0) -- (0:1.8) node[right] {\small $F$};
        \draw[->, very thick, red] (0,0) -- ($(60:1.8)+(0:1.8)$);
        \draw[dashed] (60:1.8) -- ($(60:1.8)+(0:1.8)$);
        \draw[dashed] (0:1.8) -- ($(60:1.8)+(0:1.8)$);
        \draw (0.5,0) arc (0:60:0.5);
        \node at (0.85,0.4) {\scriptsize $60^\circ$};
    \end{scope}
    % 120°
    \begin{scope}[xshift=5cm]
        \draw[->, thick] (0,0) -- (120:1.6) node[above left] {\small $F$};
        \draw[->, thick] (0,0) -- (0:1.6) node[right] {\small $F$};
        \draw[->, very thick, red] (0,0) -- ($(120:1.6)+(0:1.6)$);
        \draw[dashed] (120:1.6) -- ($(120:1.6)+(0:1.6)$);
        \draw[dashed] (0:1.6) -- ($(120:1.6)+(0:1.6)$);
        \draw (0.5,0) arc (0:120:0.5);
        \node at (0.15,0.75) {\scriptsize $120^\circ$};
    \end{scope}
    % 90°
    \begin{scope}[xshift=9.5cm]
        \draw[->, thick] (0,0) -- (0,1.7) node[above] {\small $F$};
        \draw[->, thick] (0,0) -- (1.7,0) node[right] {\small $F$};
        \draw[->, very thick, red] (0,0) -- (1.7,1.7);
        \draw[dashed] (0,1.7) -- (1.7,1.7);
        \draw[dashed] (1.7,0) -- (1.7,1.7);
        \draw (0.4,0) arc (0:90:0.4);
    \end{scope}
\end{tikzpicture}
\end{center}
